\documentclass[12pt]{article}

\textwidth 16cm \textheight 23cm \evensidemargin 0cm
\oddsidemargin 0cm \topmargin -2cm
\parindent 0pt
\parskip \medskipamount


\usepackage[dutch]{babel}
\usepackage{amssymb}
\usepackage{amsmath}
\usepackage[utf8,latin1]{inputenc}
\usepackage[normalem]{ulem} % strikethrough normal text with \sout{text}
\usepackage{cancel} % strikethrough in math mode with \cancel{text}
%\usepackage{gensymb}

\usepackage{subfig}
\usepackage{wrapfig}
\usepackage{graphicx}
\usepackage{graphics}
\usepackage{latexsym}
%\usepackage[latin1]{inputenc} %letters met accenten mogen in broncode gebruikt worden
\usepackage{amssymb,amsmath,amsthm} %allerlei wiskundige symbolen => http://www.ams.org/tex/amslatex.html voor read-me's
%\usepackage[small,bf,hang]{caption}
\usepackage[table]{xcolor}
\usepackage{pgf,tikz}
\usetikzlibrary{arrows}

\usepackage[pdftex]{hyperref} 
\hypersetup{setpagesize=false}
%\hypersetup{bookmarksopen}
\hypersetup{bookmarksnumbered} %voor de bookmarks in pdf de nrs zetten ook
\hypersetup{pdfborder=0.1pt} %de rand rond de links smaller maken
\hypersetup{linktocpage=true} %titels of pgnrs omkaderen in toc
\hypersetup{colorlinks=true} %kaders rond links of tekst in kleur, kaders gaan weg bij afprinten, kleur blijft


\numberwithin{equation}{section} %om de equations labels te geven die met section meegaan, dus niet gwn (1) (2) ... maar (1.1) (1.2) ...

%iets voor modulonotatie
\makeatletter
\def\imod#1{\allowbreak\mkern10mu({\operator@font mod}\,\,#1)}
\makeatother

%Getallenverzamelingen
\newcommand{\N}{\mathbb{N}}
\newcommand{\Z}{\mathbb{Z}}
\newcommand{\Q}{\mathbb{Q}}
\newcommand{\R}{\mathbb{R}}
\newcommand{\C}{\mathbb{C}}
\newcommand{\mO}{\mathcal{O}}
%<= en >= mooier maken
\renewcommand{\leq}{\leqslant}
\renewcommand{\geq}{\geqslant}

%speciale wiskunde letters: \mathcal{} of \mathfrak{}

%Nieuwe omgevingen voor definities en stellingen
\theoremstyle{plain}  \newtheorem{stel}{Stelling}[section]
\theoremstyle{plain}  \newtheorem{lemma}[stel]{Lemma}
\theoremstyle{plain}  \newtheorem{gevolg}[stel]{Gevolg}
\theoremstyle{definition}  \newtheorem{defi}[stel]{Definitie}
\theoremstyle{definition} \newtheorem{vb}[stel]{Voorbeeld}
\theoremstyle{definition} \newtheorem{vbn}[stel]{Voorbeelden}
\theoremstyle{definition}  \newtheorem{opm}[stel]{Opmerking}
\theoremstyle{plain}  \newtheorem{opdracht}{Opdracht}

\hyphenation{e-ner-zijds}
\hyphenation{ka-nons-ko-gel}
\hyphenation{schrij-ven}
\hyphenation{ka-nons-ko-gels}
\hyphenation{ont-stond}
\hyphenation{kwa-dra-tisch}

\usepackage{color}
\newcommand{\todo}[1]{\textcolor{red}{\##1\#}}
\newcommand{\question}[1]{\textcolor{blue}{\##1\#}}

% Wissel tussen versie van leerkracht en versie voor leerlingen
% gebruik: \teacher{Dit is een notitie voor de leerkracht}
% eerste parameter is de notitie
\newcommand{\teacher}[1]{\textcolor{gray}{\emph{#1}}}
%\newcommand{\teacher}[1]{}

% Wissel tussen versie van leerkracht en versie voor leerling door antwoorden te verbergen
% gebruik: \answer[2cm]{Dit is een antwoord}
% tussen vierkante haakjes is optioneel de lengte die voorzien moet worden (default = 0cm),
% de parameter is het antwoord
\newcommand{\answer}[2][0cm]{{\textcolor{gray}{\emph{#2}}}}
%\newcommand{\answer}[2][0cm]{\vspace*{#1}}

\graphicspath{{../cursus/figuren/}}

%\newtheorem{definitie}{Definitie}

\newcommand{\degree}{\ensuremath{^\circ}}

\begin{document}

\thispagestyle{empty}

\begin{center}
 \begin{minipage}{7cm}
\begin{center}
%\aula 
%B


\includegraphics[width=5cm]{ruglogo.pdf}


 {\large
Universiteit Gent \\
Faculteit  Wetenschappen \\
Vakgroep  Wiskunde
} \end{center}\end{minipage}

\vspace{3cm}

 {\Huge\bf 
  Vullen van Vlakken en Volumes:\\ Administratie}
  
\baselineskip=12pt
\vspace{1.5cm}

{\large\bf  Lien Lambert }\\
\vspace{0.5cm}
{\large \bf Pieter Pareit }\\
\vspace{0.5cm}
{\large \bf Jordy Vanpoucke}\\

\baselineskip=12pt

\vspace{2cm}

 
\bigskip

{\large Academiejaar 2011-2012}

%\vspace{5cm}
%\baselineskip=1.2cm



\end{center}

\vfill

\hspace*{\fill}
\begin{minipage}{8cm}
\noindent   Project Vakdidactiek Wiskunde II 
 \end{minipage}


\bigskip

\hspace*{\fill}
\begin{minipage}{8cm}
\noindent  Prof. Dr.  F. De Clerck\\
Praktijkassistent G. Maes
 \end{minipage}

%\newpage
%\tableofcontents 
%\newpage

\section{Doelgroep}

Leerlingen uit het tweede leerjaar van de tweede graad ASO. (Uiteraard is dit pakket dan ook toegankelijk voor leerlingen uit de derde graad, maar misschien iets te eenvoudig.)

\section{Voorkennis}
In dit lessenpakket over het vullen van vlakken en volumes steunen we expliciet op de geziene leerstof uit
\begin{itemize}
\item de getallenleer en algebra: De leerlingen
\begin{itemize}
\item (19) kunnen vergelijkingen van de eerste en de tweede graad in \'{e}\'{e}n onbekende oplossen,
\item (20) kunnen ongelijkheden van de eerste en de tweede graad in \'{e}\'{e}n onbekende oplossen,
\item (21) lossen problemen op die kunnen vertaald worden naar:
\begin{itemize}
\item een vergelijking van de eerste en de tweede graad in \'{e}\'{e}n onbekende.
\item een ongelijkheid van de eerste en de tweede graad in \'{e}\'{e}n onbekende.
\end{itemize}
\end{itemize}
\item de re\"{e}le functies: De leerlingen
\begin{itemize}
\item (28) kunnen stelsels van twee vergelijkingen van de eerste graad met twee onbekenden algebra�sch oplossen en de oplossing grafisch interpreteren,
\item (29) kunnen problemen oplossen die te vertalen zijn naar stelsels van twee vergelijkingen van de eerste graad met twee onbekenden,
\item (31) lossen problemen op die kunnen beschreven worden met eerste- en tweedegraadsfuncties.
\end{itemize}
\item de meetkunde: De leerlingen
\begin{itemize}
\item (34) verklaren gelijkvormigheid van figuren met behulp van schaal en congruentie,
\item (35) gebruiken de gelijkvormigheid van driehoeken en de stelling van Thales om de lengte van lijnstukken te berekenen,
\item (39) kunnen problemen met zijden en hoeken van driehoeken uit de technische wereld oplossen door een effici�nte keuze te maken uit:
\begin{itemize}
\item de stelling van Thales
\item de stelling van Pythagoras
\item goniometrische getallen
\end{itemize}
\item (41) lossen eenvoudige problemen i.v.m. ruimtelijke situaties op door gebruik te maken van eigenschappen van vlakke figuren,
\item (43) kunnen de inhoud van sommige ruimtelijke objecten benaderend berekenen door ze op te splitsen in of aan te vullen tot gekende lichamen,
\item (44) kunnen effecten van schaalverandering op inhoud en oppervlakte berekenen.
\end{itemize}
\end{itemize}


\section{Globale motivatie}
Voor deze lessenreeks zijn we vertrokken vanuit het bolstapelprobleem van Kepler. Het viel ons op hoe men in het dagelijkse leven bezig is met het zo optimaal mogelijk stapel van bollen, maar ook van kubussen en dergelijke. In deze lessenreeks wordt dit opgebouwd door te starten met het tweedimensionale probleem, namelijk de vlakvulling. In deze lessenreeks is het de bedoeling dat de leerlingen zelf gaan experimenteren, en hun bevindingen ook hard maken via het bewijzen of uitrekenen ervan.

\subsection{Vakoverschrijdende eindtermen}
In het lessenpakket komen de volgende vakoverschrijdende eindtermen aan bod.
\subsubsection{Gemeenschappelijke stam}
De leerlingen	 
\begin{itemize}
\item	(2) kunnen originele idee\"{e}n en oplossingen ontwikkelen en uitvoeren,
\item (4)	blijven, ondanks moeilijkheden, een doel nastreven,
\item (6)	kunnen schoonheid ervaren, 
\item (7)	kunnen schoonheid cre\"{e}ren,
\item (8)	benutten leerkansen in diverse situaties,
\item (10) engageren zich spontaan,
\item (11) kunnen gegevens, handelwijzen en redeneringen ter discussie stellen a.d.h. van relevante criteria,
\item (12) zijn bekwaam om alternatieven af te wegen en een bewuste keuze te maken,
\item (13) kunnen onderwerpen benaderen vanuit verschillende invalshoeken,
\item (19) dragen actief bij tot het realiseren van gemeenschappelijke doelen.
\end{itemize}
\subsubsection{Leren leren}
De leerlingen
\begin{itemize}
\item (4)	kunnen verwerkte informatie vakoverstijgend en in verschillende situaties functioneel toepassen,
\item (5)	kunnen informatie samenvatten,
\item (6)	kunnen op basis van hypothesen en verwachtingen mogelijke oplossingswijzen realistisch inschatten en uitvoeren,
\item (7)	evalueren de gekozen oplossingswijze en de oplossing en gaan eventueel op zoek naar een alternatief.
\end{itemize}
\subsubsection{Contexten}
We maken enkel gebruik van eindtermen uit context 7, `socioculturele samenleving'.\\
De leerlingen
\begin{itemize}
\item (6) gaan actief om met de cultuur en kunst die hen omringen,
\item (7)	illustreren de wederzijdse be�nvloeding van kunst, cultuur en techniek, politiek, economie, wetenschappen en levensbeschouwing.
\end{itemize}
\subsection{Vakgebonden eindtermen}
We maken gebruik van de volgende vakgebonden eindtermen.\\
De leerlingen kunnen
\begin{itemize}
\item (1)	wiskundetaal begrijpen en gebruiken,
\item (2)	wiskundige informatie analyseren, schematiseren en structureren,
\item (3)	eenvoudig mathematiseerbare problemen ontleden (onderscheid maken tussen gegevens en gevraagde, de relevantie van de gegevens nagaan en verbanden leggen ertussen) en vertalen naar een passende wiskundige context,
\item (4)	wiskundige problemen planmatig aanpakken (door eventueel hi\"{e}rarchisch op te splitsen in deelproblemen),
\item (5)	bij het oplossen van wiskundige problemen kritisch reflecteren over het oplossingsproces en het eindresultaat,
\item (8)	voorbeelden geven van de rol van de wiskunde in de kunst,
\item (9)	kennis, inzicht en vaardigheden die ze verwerven in wiskunde bij het verkennen, vertolken en verklaren van problemen uit de realiteit gebruiken.
\end{itemize}
De leerlingen
\begin{itemize}
\item (11*)	leggen een zin voor nauwkeurigheid aan de dag bij het hanteren en het toepassen van wiskunde,
\item (13*)	zijn gericht op samenwerking om de eigen mogelijkheden te vergroten.
\end{itemize}

\section{Planning, lesdoelstellingen, werkvormen en didactische wenken}
\subsection{Les 1: Inleiding}

\subsection{Les 2: Vullen van vierkanten met kleinere vierkanten en zo tot Perfecte Vierkanten komen}
\subsubsection{Lesdoelstellingen}

\begin{itemize}
  \item De leerlingen kunnen een een vierkant met kleinere vierkanten vullen van de zelfde grootte, om zo een algemene formule die de maximale effici�ntie bepaald, te vinden.
  \item De leerlingen weten wat een perfect vierkant van orde $n$ is.
  \item De leerlingen kunnen een ondergrens zoeken voor de orde van het perfect vierkant door redeneren.
  \item De leerlingen zien in dat een perfect vierkant dat gedraaid of gespiegeld wordt geen ander perfect vierkant oplevert. Hiervoor leren de leerlingen het begrip isomorfisme.
  \item De leerlingen beseffen dat het vinden van ��n perfect vierkant leidt tot het vinden van een oneindig aantal perfecte vierkanten. 
  \item De leerlingen kunnen een bewijs uit het ongerijmde opstellen om een eigenschap te bewijzen over perfecte vierkanten.
  \item De leerlingen kunnen hun resultaten en definities in verband met perfecte vierkanten uitbreiden naar ��n en drie dimensies.
\end{itemize}

\subsubsection{Werkvormen}

\begin{itemize}
  \item Onderwijsleergesprek.
  \item Individueel werken.
  \item Leerlingen laten puzzelen.
\end{itemize}

\subsubsection{Didactische wenken}

Het is de bedoeling dat de leerkracht samen met de leerlingen een theorie rond perfecte vierkanten ontwikkelt.

We beginnen samen met de leerlingen aan het vullen van een vierkant met gelijke vierkanten. Hiervoor berekenen we dan de effici�ntie, dus hoeveel van het vierkant wordt bedekt door kleinere vierkanten. Hiervoor kunnen we snel een algemene formule vinden.

Dan wensen we het probleem moeilijker maken. Hiervoor eisen we dat een vierkant volledig gevuld wordt met kleinere vierkanten die allemaal een verschillende grootte hebben. Omdat dit een moeilijk probleem is, bijvoorbeeld naar het perfect vierkant met de laagste orde is 60 jaar gezocht, geven we in deze cursus zelf een voorbeeld en een beetje geschiedenis, wat de leerkracht aanbrengt.

De leerlingen die toch zelf actief aan de slag willen gaan, en een perfect vierkant puzzelen, geven we deze kans. Voor hen werden in de appendix de bouwplannen voorzien voor alle perfecte vierkanten van orde 22. Een perfect vierkant vinden is heel moeilijk, maar wanneer je reeds weet wat de groottes zijn van de kleinere vierkantjes wordt het toch haalbaar.

Voorgaande is optioneel en kan gerust door de leerlingen die graag puzzelen thuis gedaan worden. Een competitie is ook mogelijk. Hierna gaan we door. We kunnen aantonen dat de orde minimaal een bepaalde grootte moet zijn. Voor een lage orde lukt dit. We nemen daarna ��n bepaald perfect vierkant en tonen dan aan dat vanuit dit perfecte vierkant nog een ander perfect vierkant geconstrueerd kan worden.

Als laatste bewijzen we een bepaalde eigenschap over een algemeen perfect vierkant. Deze eigenschap wordt dan gebruik om aan te tonen dat een perfect vierkant niet uitgebreid kan worden tot een perfecte kubus. Deze eigenschappen worden bewezen uit het ongerijmde, maar kunnen heel visueel aangepakt worden. Hiervoor moeten dan wel een paar vierkanten van verschillende grootte uit een stuk karton geknipt worden. 


\subsection{Les 3: Puzzelen met vierkanten en rechthoeken om vergelijkingen op te lossen}

\subsection{Les 4: Kennismaking met tweedimensionale stapelproblemen en hiermee experimenteren}

\subsubsection{Lesdoelstellingen}
\begin{itemize}
\item De leerlingen beseffen dat je een tafel niet volledig kan bedekken met cirkels
en dat we daarom op zoek kunnen gaan naar de optimale configuratie, namelijk 
het hexagonaal rooster.

\item De leerlingen weten wat een voronoicel is en kunnen hiermee de effici�ntie
 van een rooster berekenen.

\item De leerlingen kunnnen de oppervlakte van een driehoek, rechthoek, zeshoek en cirkel
berekenen.

\item De leerlingen kunnen a.d.h.v. een geogebra-applet of door te experimenteren met jetons of munten de effici�ntie van dienbladen bepalen.

\item De leerlingen kunnen de effici�ntie voor eenvoudige voorbeelden uitrekenen.

\item De leerlingen ontwikkelen meetkundig inzicht.

\end{itemize}

\subsubsection{Werkvormen}
\begin{itemize}
\item Onderwijsleergesprek
\item groepswerk
\item leerlingen laten experimenteren
\end{itemize}
\subsubsection{Didactische wenken}
Op weg naar he bolstapelprobleem van Kepler houden we eerst halt bij het tweedimensionael equivalent. Hiervoor trekken we drie lessen uit.

In het eerste deel van de eerste les van \textquoteleft tweedimensionale stapelproblemen' werken de leerlingen samen met de leerkracht aan de hand van een bundeltje. Het is aangewezen dat de leerlingen al in kleine groepjes samen zitten, waarbij de leerkracht echter eerst nog vooraan in de klas de leider is van het onderwijsleergesprek. In het bundeltje staan vragen en opdrachten. Het is de bedoeling dat de leerkracht de leerlingen aanzet om hierover na te denken of effectief de opdracht uit te voeren. Dit kan gaan van een kleine berekening tot effectief experimenteren met munten. 

Als we bij het dienbladenprobleem zijn aanbeland is het effectief aan de leerlingen om aan de slag te gaan. Ze moeten voor verschillende aantallen munten en verschillende vormen van dienbladen (cirkelvormig, driehoekig, zeshoekig, rechthoekig) de effici�ntie berkenen. Het is de bedoeling dat ze dit experimenteel uitvoeren, a.d.h.v. Geogebra-applets en ook aan de hand van echt materiaal (jetons en karton). 
Een deel van de groepjes kan beginnen aan de computers, een ander deel kan starten met het werken met het andere materiaal. Het is echter ook de bedoeling om de leerlingen sommige configuraties ook echt te laten tekenen en de effici�ntie te laten berekenen.
De leerkracht kan de leerlingen zeggen dat ze dit voor een bepaald aantal moeten doen of voor enkele specifieke gevallen.
Je kan de klas ook opsplitsen, zodanig dat niet iedereen alle effici�nties moet berekenen, maar dat je met de ganse klas weldegelijk komt tot een volledig ingevulde tabel.


\subsection{Les 5: Het dienbladenprobleem: experimenteren en effectief berekenen}
\subsubsection{Lesdoelstellingen}

\begin{itemize}
\item De leerlingen kunnen a.d.h.v. een geogebra-applet of door te experimenteren met jetons of munten
de effici�ntie van dienbladen bepalen.
\item De leerlingen kunnen functioneel gebruik maken van ict.

\item De leerlingen kunnen de effici�ntie voor eenvoudige voorbeelden uitrekenen.

\item De leerlingen kunnen vanuit een figuur meetkundig inzicht verwerven en hierdoor 
problemen vereenvoudigen.

\item De leerlingen kunnen de meetkundige interpretatie van de goniometrische formules functioneel gebruiken, i.h.b. voor projecties.

\end{itemize}

\subsubsection{Werkvormen}
\begin{itemize}
\item Onderwijsleergesprek
\item groepswerk
\item leerlingen laten experimenteren
\end{itemize}

\subsubsection{Didactische wenken}

In deze les gaan de leerlingen verder met experimenteren en met het berekenen van de effici�ntie van sommige gevallen. De leerkracht blijft in de klas rondlopen en helpt waar nodig.
Als de klas was opgedeeld in een helft aan de computer en een helft met het andere materiaal, is het aangewezen om nu te wisselen.
In het tweede deel van deze les is het al handig om de leerlingen samen te brengen om wat ze hebben gevonden samen te leggen en samen met de leerlingen de effici�nties van de configuraties met 1, 2 of 3 munten te overlopen.


\subsection{Les 6: Het dienbladenprobleem: berekenen, bewijzen en besluiten}
\subsubsection{Lesdoelstellingen}
\begin{itemize}
\item De leerlingen beseffen dat men uit experimenteren belangrijke informatie
kan halen maar beseffen ook het belang van bewijzen.

\item De leerlingen kunnen aan de hand van een stappenplan het bewijs zelf verwerven.

\item De leerlingen kunnen de effici�ntie van iets ingewikkeldere configuraties berekenen.

\item De leerlingen kunnen vorige berekeningen gebruiken om ingewikkelde configuraties te vereenvoudigen.

\item De leerlingen kunnen uit experimentele resultaten conclusies trekken over de meest
effici�nte dienbladen, en kunnen ook (mogelijke) verklaringen geven

\end{itemize}



\subsubsection{Werkvormen}
\begin{itemize}
\item Onderwijsleergesprek
\item groepswerk
\item leerlingen laten experimenteren
\end{itemize}

\subsubsection{Didactische wenken}
Deze les is vooral bedoeld voor de leerlingen om hen te doen inzien dat ze
het ook belangrijk is om niet steeds zomaar op de intu�tie in te gaan. Het uitproberen met Geogebra is leuk, het experimenteren met figuren ook. Je ziet dan snel in `op het gevoel' dat je niet effici�enter kan, maar het is belangrijk om dit ook te kunnen aantonen. We zullen het geval behandelen van een driehoekig dienblad met vier munten. In ons geval is onze intu�tie weldegelijk correct volgens het bewijs, nu zijn we dus zeker. De leerlingen kunnen dit bewijs zelf doorlopen met deze cursus en dit begrijpen via een figuur die ze zelf knutselen.

Op het einde van de les kan de leerkracht nog eens terugkomen op enkele andere specifieke gevallen voor het berekenen van de configuratie.

Belangrijk is om op het einde de volledige tabel, die in de vorige les normaalgezien al was samengelegd, te overlopen en er de belangrijkste zaken uit te halen over wat nu de effici�ntste dienbladen zijn en als dit klopt met de intu�tie en met wat in werkelijkheid wordt gebruikt.



\subsection{Les 7: Kanonskogels stapelen}

\subsubsection{Lesdoelstellingen}

\begin{itemize}
  \item De leerlingen weten wat het bolstapelprobleem inhoudt.
	\item De leerlingen kunnen het aantal kanonskogels in een (regelmatige) piramide met vierkant en driehoekig grondvlak met $n$ lagen experimenteel berekenen.
	\item De leerlingen kunnen vanuit een gegeven rij het recursief en/of expliciet voorschrift geven.
	\item De leerlingen kunnen door het doordacht experimenteren bepaalde vragen over combinaties van piramides onderzoeken.
	\item De leeringen kunnen een bewijs geven voor de som van de eerste $n$ kwadraten.
	\item De leerlingen kunnen op een analoge wijze het bewijs geven voor de som van de eerst $n$ driehoeksgetallen.
\end{itemize}

\subsubsection{Werkvormen}
\begin{itemize}
\item Onderwijsleergesprek
\item In groepjes werken en experimenteren
\end{itemize}

\subsubsection{Didactische wenken}
Na een korte inleiding van de leerkracht, geeft hij/zij de leerlingen de opdracht om op zoek te gaan naar het aantal kanonskogels in een (regelmatige) piramide met vierkant en driehoekig grondvlak met $n$ lagen. Om dit uit te testen krijgen de leerlingen bijvoorbeeld knikkers of andere bolvormige voorwerpen. Ze kunnen eerst dit eenvoudig uittesten voor enkele kleine gevallen. Hieruit is het dan de bedoeling dat ze een formule vinden die afhangt van $n$, het aantal lagen van de piramide. De leerkracht loopt ondertussen rond en helpt waar mogelijk. Het is zeker de bedoeling dat de leerlingen gevonden hebben dat het de som is van de eerste $n$ kwadraten en eerste $n$ driehoeksgetallen.
Ook krijgen de leerlingen een aantal andere vragen mee die ze in kleine groepjes moeten gaan uitproberen. Uiteraard hoeven de leerlingen niet constant met de knikkers te werken. Eenmaal ze het systeem doorhebben, kunnen ze met de rijen op papier werken en de vragen oplossen.

Op het einde van de les is het de bedoeling dat de leerkracht met de leerling kijkt of ze de formule hebben gevonden en dat de leerlingen aan elkaar kunnen uitleggen hoe ze daaraan gekomen zijn. Daarna is het bedoeling dat de leerkracht met de leerlingen de expliciete formule geeft voor de eerste $n$ kwadraten. Op analoge wijze kunnen de leerlingen nu de formule van de eerste $n$ driehoeksgetallen vinden. Dit kan ook eventueel als extra opgave gegeven worden, die niet noodzakelijk in de klas moet behandeld worden, afhankelijk van de tijd.

\subsection{Les 8: Vervolg kanonskogels en het bolstapelprobleem van Kepler}

\subsubsection{Lesdoelstellingen}

\begin{itemize}
\item De leerlingen kunnen door het doordacht experimenteren bepaalde vragen over combinaties van piramides onderzoeken.
\item De leerlingen kunnen vanuit een voorbeeld komen tot een algemene formule en dit ook bewijzen, al dan niet met behulp van de leerkracht.
	\item De leerlingen weten wat het bolstapelprobleem inhoudt.
	\item De leerlingen beseffen dat het vinden van een bewijs soms vele jaren en generaties kan duren met vallen en opstaan.
	\item De leerlingen weten wat een \textquoteleft kissing number' is.
\end{itemize}

\subsubsection{Werkvormen}
\begin{itemize}
\item Onderwijsleergesprek
\item quiz
\end{itemize}

\subsubsection{Didactische wenken}
We gaan in deze les nog even verder met waar we de vorige les nog niet aan toe komen, namelijk het combineren van stapels van kanonskogels. De leerkracht overloopt samen met de leerlingen hun bevinden welke combinaties er mogelijk zijn, welke zeker niet, waarom niet enzovoort. Uiteindelijk is het de bedoeling dat voor de vragen waar we een positief antwoord op kunnen geven, we vertrekken vanuit een voorbeeld om hieruit een algemene formule voor te vinden en dit effectief ook te bewijzen met de verkregen formules van de vorige les.

Voor het tweede deel van de les is het de bedoeling dat we op een ludieke manier de evolutie van het bolstapelprobleem van Kepler schetsen. We doen dit aan de hand van een quiz. De leerlingen worden verdeeld in groepjes van ongeveer 4 personen en krijgen allemaal bordjes of briefjes met A, B, C en D op. Het is namelijk een quiz met meerkeuzevragen. Sommige antwoorden zouden de leerlingen tijdens de loop van deze lessenreeks al moeten gehoord hebben, andere vragen zijn eerder \textquoteleft gokvragen'. De leerkracht houdt ook de stand bij, voor elk goed antwoord krijgen de leerlingen ��n punt. Er zijn 8 vragen en bij elke vraag moeten de leerlingen op het teken van de leerkracht hun bordjes omhoog steken.
Deze quiz is echter maar een werkvorm om het verhaal van het bolstapelprobleem te kunnen aanbrengen.
Het is de bedoeling dat de leerlingen beseffen wat voor een evolutie achter bepaalde wiskundige problemen zitten en dat dit zoeken met vallen en opstaan gebeurt.

\end{document}
